\documentclass[11pt,a4paper]{article}
\usepackage{amsmath, amssymb, fullpage, mathrsfs, bm, pgf, tikz, fullpage, amsthm}
\usepackage{mathrsfs, float}
\usetikzlibrary{arrows}
\setlength{\textheight}{10in}
%\setlength{\topmargin}{0in}
\setlength{\topmargin}{-0.5in}
\setlength{\parskip}{0.1in}
\setlength{\parindent}{0in}


\begin{document}
		\newcommand{\la}{\leftarrow}
	\newcommand{\lra}{\leftrightarrow}
	\newcommand{\bbN}{{\mathbb N}}
	\newcommand{\bbZ}{{\mathbb Z}}
	\newcommand{\bbQ}{{\mathbb Q}}
	\newcommand{\bbR}{{\mathbb R}}
	\newcommand{\bbC}{{\mathbb C}}
	\newcommand{\bbH}{{\mathbb H}}
	\newcommand{\dfeq}{\stackrel{\mathrm{def}}{=}}
	\newcommand{\ra}{\rightarrow}
	\newcommand{\Span}{\mathrm{span}}
	\newcommand{\scrP}{\mathscr{P}}
	\newcommand{\rank}{\mathrm{rank}}
	\newcommand{\nullity}{\mathrm{nullity}}
	\newcommand{\Col}{\mathrm{Col}}
	\newcommand{\Row}{\mathrm{Row}}
	\newcommand{\tr}{\mathrm{tr}}
	\newcommand{\ol}{\overline}
	\newcommand{\norm}[1]{||#1||}
	\newcommand{\doubleline}[1]{\underline{\underline{#1}}}
	\newcommand{\elemop}[1]{\stackrel{#1}{\longrightarrow}}
	\newcommand{\Ind}{\mathrm{Ind}}
	\newcommand{\Res}{\mathrm{Res}}
	\newcommand{\End}{\mathrm{End}}
	\newcommand{\cl}{\mathrm{cl}}
	\newcommand{\code}[1]{\texttt{#1}}
	\newcommand\tab[1][0.5cm]{\hspace*{#1}}
	\newcommand{\<}{\langle}
	\renewcommand{\>}{\rangle}
	\newcommand{\qubits}[1]{|{#1}\rangle}
	\newcommand{\powset}{\mathcal{P}}
	\newcommand{\dsum}{\displaystyle\sum}
	\newcommand{\dprod}{\displaystyle\prod}
	
	\newtheorem{lemma}{Lemma}
	
	\title{Solution (sketch) to IMO 2025 Shortlist}
	\date{}
	\author{Anzo Teh}
	\maketitle
	
	\section*{Algebra}
	\begin{enumerate}
		\item [A1.] 
		Quadratic solitaire is a single-player game. To start the game, the player chooses two distinct nonzero integers $a$ and $b$ and writes the equation $x^2+ax+b=0$ on a blackboard. On each turn, if the equation currently written on the blackboard has two distinct nonzero integer solutions $x=u$ and $x=v$, then the player erases the equation and replaces it by one of the equations $x^2+ux+v=0$ and $x^2+vx+u=0$ of their choosing. Otherwise, the game ends. Determine all initial choices of $a$ and $b$ such that quadratic solitaire can be played forever.
		
		\textbf{Answer.}$(1, -2)$ where $x^2+x-2=(x-1)(x+2)$ gives the same polynomial. 
		
		\textbf{Solution.}
		Now note that $uv = b$ so $|b|\ge \max\{|u|, |v|\}$ (as each $|u|, |v|\ge 1$)
		so $|b|$ is nonincreasing, hence needs to stabilize. 
		Note that $|v|=|b|$ means $|u|=1$, hence eventually we end up with $x^2\pm x + v$ with one of $\{1, -1\}$ as root. 
		Since $v\neq 0$, $x^2\pm x\in\{-2, 0, 2\}$ when $x=\pm 1$, 
		so $|v|=2$. $x^2\pm x+2$ has no real root as discriminant is $-7$, hence $v=-2$. 
		Now $x^2-x-2=(x+1)(x-2)$ has root $\{-1, 2\}$, 
		but arrangement into $x^2-x+2$ has no real root. 
		Thus the only `terminal' polynomial is $x^2+x-2$. 
		
		\item [A2.]
		The sunshine cost of a sequence $a_1$, $a_2$, $\ldots$, $a_{100}$ of integers is the largest possible value of
		\[|(a_1+a_2+\cdots+a_i)-a_j|\]as $i$ and $j$ vary over all integers $1$, $2$, $\ldots$ , $100$.
		
		
		Determine the smallest possible sunshine cost over all sequences $a_1$, $a_2$, $\ldots$, $a_{100}$ of pairwise distinct integers.
		
		\textbf{Answer.} 75, or $\lceil \frac{3(n-1)}{4}\rceil$ by replacing 100 with general $n$ that's multiple of 4. 
		
		\textbf{Solution.}
		\textit{Lower bound}: we can't choose $a_1$ as the maximizer of $|a_i|$; otherwise either $a_1$ is the max or min and $|a_1-a_i|\ge n - 1$ for some $i$. 
		
		Let $\max a_j-a_i=M$ satisfied by $(j_0, i_0)$ and $\max |a_i|=C$; w.l.o.g. let $a_k = C$, $k\ge 2$. 
		clearly can't have $|a_1|=C$. 
		Let the sum $(a_1+\cdots + a_i)-a_j=C_{i, j}$. 
		Now 
		\[
		C_{k, i} - C_{k - 1, j} = a_k + a_j - a_i = C + M
		\]
		so the answer is at least $\frac{C+M}{2}$. 
		Moreover $M\ge n - 1, C\ge \lceil \frac{n - 1}{2}\rceil$, giving the LB. 
		
		Upper bound: we'll only show the case $n=4k$ as requested in the problem, and the desired bound is $3k$. We'll make sure all $|a_i|\le 2k$; all $|a_1+\cdots + a_i|\le k$. We arrange as 
		\[
		-k, 2k, -2k, 0, 1, 2k-1, -1, -(2k-1), \cdots, k-1, k+1, -(k-1), -(k+1)
		\]
		
		\item [A3.] (IMO 5)
		Alice and Bazza are playing the inekoalaty game, a two‑player game whose rules depend on a positive real number $\lambda$ which is known to both players. On the $n$th turn of the game (starting with $n=1$) the following happens:
		If $n$ is odd, Alice chooses a nonnegative real number $x_n$ such that
		\[
		x_1 + x_2 + \cdots + x_n \le \lambda n.
		\]
		If $n$ is even, Bazza chooses a nonnegative real number $x_n$ such that
		\[
		x_1^2 + x_2^2 + \cdots + x_n^2 \le n.
		\]
		If a player cannot choose a suitable $x_n$, the game ends and the other player wins. If the game goes on forever, neither player wins. All chosen numbers are known to both players.
		
		Determine all values of $\lambda$ for which Alice has a winning strategy and all those for which Bazza has a winning strategy.
		
		\textbf{Answer.} $A$ wins when $\lambda > \frac{\sqrt{2}}{2}$; 
		$B$ wins when $\lambda < \frac{\sqrt{2}}{2}$. 
		
		\textbf{Solution.}
		Let $C = \frac{\sqrt{2}}{2}$. 
		
		$\lambda > C$: 
		Coral outputs 0 at all times; 
		Joey's outputs will make sure $\sum x_{n}\le Cn$. 
		Wait till $n$ big enough (and even) then pick $x_{n+1}=(\lambda - C)n$ (then $\sum x_{n}^2 \ge (\lambda - C)n^2 > n$), Joey loses. 
		
		$\lambda < C$: 
		Joey always tries to `max out', 
		i.e. $x_{2k+1}^2+x_{2k+2}^2 = 2$. 
		Note that $x_{2k+1}+x_{2k+2}\ge \sqrt{2}$ so Coral loses at one point. 
		
		$\lambda = C$: 
		Coral sustains by outputting 0; 
		Joey sustains by making sure $x_{2k+1}^2+x_{2k+2}^2 = 2$. 
		
		\item [A4.]
		Let $\mathbb{Z}_{\ge0}$ be the set of all nonnegative integers. Let $f:\mathbb{Z}_{\ge0}\rightarrow\mathbb{Z}_{\ge0}$ be an unbounded function such that, if $m$ and $n$ are nonnegative integers satisfying
		\[ f(m+n)=\max\{f(0),f(1),\dots,f(m+n)\}, \]then
		\[ f(m+n)=f(m)+f(n). \]Prove that there exist positive integers $A, B, C$ and $D$ such that for all nonnegative integers $n$,
		\[ f(An+B)=Cn+D. \]
		
		\textbf{Solution.} 
		Denote peak as $k\ge 1$ satisfying $f(k)\ge f(\ell)$ for all $\ell\le k$. Let $C = \min \{f(m): m\ge 1\}$. There are infinitely many peaks since $f$ is unbouned. We show the following: 
		
		\textit{$C > 0$.} Indeed, if $f(m) = 0$ for $m\ge 1$ then for any peak $n\ge m$ we have $f(n)=f(n-m)$ and so $f$ is bounded. 
		
		\textit{Exactly one $m$ has $f(m)=c$.} Let $m_1 < m_2$ be such that $f(m_1)=f(m_2)=C$. 
		For any peak $n > m_2$, we have $f(n)=f(n-m_1)+C=f(n-m_2)+C$ and note that both $f(n-m_1), f(n-m_2)$ are $\max \{f(0), \cdots, f(n-1)\}$. 
		Since $m_1<m_2$ we have $f(n-m_1)=f(n-m_2)+f(m_1-m_2)$ and $f(m_1-m_2)=0 < C$, which is contradiction.
		
		\textit{If $n$ is a peak, $f(A)=C$, and $n>A$,  then $n - A$ is a peak.}
		For any $m < n$ we have $f(n-A)=f(n)-f(A)\ge f(n)-f(m)=f(n-m)$ so $f(n - A) = \max \{f(0), \cdots,  f(n - 1)\}$. 
		
		The last point says the number of peaks among $\{kA+1, \cdots, (k+1)A\}$ is at most that of $\{(k-1)A+1, \cdots, kA\}$; choose such $k$ such that this number stabilizes and let $B\in \{kA+1, \cdots, (k+1)A\}$ be a peak. 
		Then $B, A+B, 2A+B, \cdots$ are all peaks. Therefore $f(nA+B)=nf(A)+f(B)=Cn+D$, as claimed. 
		
		\item [A6.]
		Let $S$ be a set of positive integers, possibly infinite, such that no positive integer greater than $1$ divides all elements of $S$. Determine all non-periodic infinite sequences $a_1,a_2,a_3,\dots$ of positive integers such that, for all positive integers $n$,
		\begin{itemize}
			\item $a_n\leqslant |a_{n+\ell}-\ell|$ for all $\ell$ in $S$, and
			\item $a_n=|a_{n+\ell}-\ell|$ for at least one $\ell$ in $S$.
		\end{itemize}
		
		\textbf{Answer.} $a_n = n + c$ for $c\ge 0$, 
		which obviously fulfills. 
		
		\textbf{Solution.} We rewrite the condition in the following form: \[
		a_n = \min \{|a_{n+\ell} - \ell|: \ell\in S\} \cdots (*)
		\]
		
		\begin{lemma}
			 All bounded sequences in (*) are periodic. 
		\end{lemma}
	
	    \begin{proof}
	    	Let $M$ be $\max \{a_n\}$, then $M < |a_{n+\ell} - \ell|$ whenever $\ell > 2M$. 
	    	Denote $s_{n, k} = (a_{n+1}, \cdots, a_{n+k})$; 
	    	$a_n$ is a function of $s_{n, 2M}$. 
	    	Let $h_n$ be the minimum positive integer such that $s_{n, 2M} = s_{n - h_n, 2M}$ (or $\infty$ if this number doesn't exist), 
	    	then $s_{n, 2M} = s_{n - h_n, 2M} \to a_n = a{n - 2M} \to s_{n - 1, 2M} = s_{n - h_n - 1, 2M}$ so 
	    	$h_{n-1}\le h_n$ and given that there are at most $M^{2M}$ distinct possibilities of $s_n$, $h_n\le M^{2M}$ for infinitely many $n$'s, 
	    	and therefore $h_n\le M^{2M}$ for all $n$. 
	    	Choose $h = \max \{h_n\}$, we have $a_n$ periodic with period $h$. 
	    \end{proof}
		
		Now among unbounded sequences, it suffices to show that $a_n = n + c$ for $n$ sufficiently large by (*). 
		Denote $b_n = a_n - n$ and $m = \lim_{n\to\infty}\inf\{b_n\}$. 
		; the second bullet point $a_n = |a_{n+\ell} - \ell|$ means either $b_{n+\ell} = b_n$ or $b_{n+\ell} = -a_n - n < b_n$; 
		so either way $b_{n+\ell} \le b_n$ and iteratively, we get $m\le b_n$ for all $n$. 
		On the other hand, choose a finite set $S_1$ such that $\gcd S_1=1$, 
		$n_0$ such that $a_{n_0} > \max S_1$. Then $a_{n_0+\ell}\ge a_{n_0}+\ell$ for all $\ell\in S_1$ (i.e. $b_{n_0+\ell}\ge b_{n_0})$. 
		Let $p_0$ be such that for all integers $p\ge p_0$, $p$ can be written as nonnegative integer combinations of $S_1$ (which is valid as $\gcd S_1=1$, 
		can similarly deduce $b_{n_0+p}\ge b_{n_0}$ for all $p\ge p_0$, 
		and since $a_{n_0+p} > \max S_1$, 
		this also gives $b_{n_1+p}\ge b_{n_1}$ for all $n_1 > n_0+p$ 
		(i.e. $m\ge b_{n}$ for all such $n$'s). 
		Thus combining the arguments we have $m = b_n$ for all $n > n_1$, i.e. $a_n = n + m$ for all $n > n_1$, and so the argument follows.
		
		\item [A7.]
		Let $k\geqslant 3$ be a positive integer. Prove that there exists a unique $k$-tuplet$(x_1,x_2,\dots,x_k)$ of positive real numbers such that $x_1\geqslant x_2\geqslant\cdots \geqslant x_k$, and$$\left\lfloor nx_1+\frac{1}{2}\right\rfloor+\left\lfloor nx_2+\frac{1}{2}\right\rfloor+\cdots+\left\lfloor nx_k+\frac{1}{2}\right\rfloor=n$$for every integer $n$.
		
	    \textbf{Answer}. $x_m = \frac{2^{k - m}}{2^k - 1}$. 
		
		\textbf{Solution.} 
		Call all $k$-tuples working if it satisfies the problem condition. We start with the following set up: 
		
		\begin{itemize}
			\item Picking $n$ large and using $x - 1 <\lfloor x\rfloor \le x$, we have $\sum x_i = 1$; 
			
			\item $\lfloor y + \frac 12 \rfloor + \lfloor -y + \frac 12 \rfloor$ is 1 if $y + \frac 12$ is integer and 0 otherwise. 
			Considering $n$ and $-n$ simultaneously for each $n > 0$, $nx_m$ cannot be half-integer. 
			
			\item Let $S_x$ be integers $n$ where $\lfloor nx + \frac 12\rfloor - \lfloor (n-1)x + \frac 12\rfloor = 1$; 
			$S_{x_1}, \cdots, S_{x_k}$ are pairwise disjoint but union is $\mathbb{N}$. 
			
			\item Any $x_1, x_2\ge 0$ with sum 1 and $nx_1$ not a half-integer is a working pair. 
			
			\item Denote $T_{x}=\min S_x\cap \mathbb{N}$; then $T_{x_1}=1$. 
			In addition, given that $x_1$ is such that $nx_1$ is not a half-integer, 
			$S_{1 - x_1} = \bigcup_{m=2}^k S_{x_m}$
			and in particular, $T_{x_2} = T_{1-x_1}\ge 2$. 
			
			\item Denote $G_x$ the set of possible gaps between consecutive elements in $S_x$. 
			We have $G_x = \mathbb{N}\cap [\lfloor \frac 1x\rfloor, \lceil \frac 1x\rceil]$. 
			In addition, when $x < \frac 12$, 
			we have $\frac{1}{2(t-1)} \le x \le \frac{1}{2t}$ where $t=T_x$, so $G_x\in \{2t-2, 2t-1, 2t\}$
		\end{itemize}
	    
	    We now consider the following steps. 
	    
	    \begin{lemma}
	    	$T_{x_2}=T_{1-x_1}=2$, $G_{x_2}= \{3, 4\}$ and $G_{x_1}=\{2, 3\}$. 
	    \end{lemma}
        
        \begin{proof}
        	Denote $T = T_{x_2}$ for simplicity. 
        	By inspecting any $t\in S_{x_3} > T_{x_2}$, and $t_0, t_1$ be consecutive elements in $S_{x_2}$ such that $t_0 < t < t_1$, we have $2T\ge \max G_{x_2}\ge t_1 - t_0 \ge 2\min G_{1-x_1}\ge 2(2T-2)$, so $T\ge 2T-2$ and $T=2$. 
        	In addition, these inequalities must all hold so $\max G_{x_2}=4, \min G_{1-x_1}=2$. 
        	Finally, both of them must include 3, otherwise either $x_2=\frac 14$ or $1-x_1=\frac 12$. 
        	These cases mean $2x_2$ and $x_1$ are half-integer, which is a contradiction.
        \end{proof}
        
        Next, we note that each $S_x$ partitions the real number line into intervals. We now consider the `blocks' of intervals as follows: 
        $m$ blocks are defined by $m+1$ points including both endpoints. 
        
        \begin{lemma}
        	All consecutive blocks of interval length 2 in $S_{1-x_1}$ must contain even and equal number of intervals. 
        \end{lemma}
        
        \begin{proof}
        	If $t\in S_{1-x_1}$ is an endpoint to an interval of length 3, and $t\not\in S_{x_2}$, then $S_{x_2}$ will have a gap of a least $2+3\ge 5$ which contradicts $G_{x_2}=\{3, 4\}$. 
        	Conversely, for each of the intervals of size 2 in $S_{1-x_1}$, one of them cannot belong to $x_2$. 
        	This leaves each consecutive blocks to have even number of 2's, with the alternating ones (including both endpoints) belong to $S_{x_2}$ and the other alternating ones $\cup_{i=3}^k S_{x_i}$. 
        	
        	Now, if $a\ge 2$ is the number of consecutive intervals of size 2 in $S_x$, i.e. for some $t$, 
        	$t, t+2, \cdots, t+2a\in S_{x}$, but $t-a, t+2a+2\not\in S_x$. 
        	Then 
        	\[
        	\lfloor (t+2a)x+\frac 12)\rfloor - \lfloor (t-1)x+\frac 12)\rfloor\ge a + 1
        	\qquad 
        	\lfloor (t+2a+2)x+\frac 12)\rfloor - \lfloor (t-3)x+\frac 12)\rfloor\ge a + 1
        	\]
        	The first implies $(2a+1)x > a$ and $x > \frac{a}{2a+1}$; 
        	the second implies $(2a+5)x < a+2$ so $x < \frac{a+2}{2a+5}$. 
        	Thus there cannot be intervals of $a$ blocks and $b$ blocks simultaneously if $b \ge a + 2$, proving the claim. 
        \end{proof}
		
		Consequently, if $2a$ is the consecutive block size (of interval 2) then we in fact have $1 - x_1 = \frac{2a + 1}{4a + 3}$, with 
		\[
		S_{1-x_1} = \{m(4a + 3) + 2p, p = 1, 2, \cdots, (2a+1)\}
		\qquad 
		S_{x_2} = \{m(4a + 3) + 4p+2, p = 0, 1, \cdots, a\}
		\]
		For convenience, denote $M = 4a + 3$. 
		Our goal is to show that $M = 2^m-1$. 
		We now iteratively establish the following about $S_{x_m}$ for $m=3, \cdots, k$. 
		
		\begin{lemma}
			For all $m$ with $m\le k$, $U_{m} := \cup_{m'=m}^k S_{x_{m'}} = \{Mu + 2^{m-1}c, c = 1, 2, \cdots, \lfloor \frac{M}{2^{m-1}}\rfloor\}$, 
			and in addition $U_m\cap \{1, 2, \cdots, M\}$ has odd number of elements (so $2^m \mid M + 1$). 
		\end{lemma}
	    
	    \begin{proof}
	    	We proceed by induction on $m$. 
	    	For $m = 3$ the previous discussion already gave us $U_3 = \{4, 8, \cdots, M - 3\}\pmod{M}$, 
	    	and $4 \mid M + 1$. 
	    	Our discussion will first assume this form of $U_m$ and that $2^{m-1}\mid M + 1$, show that $U_m$ has odd number of elements, $2^m\mid M + 1$, 
	    	and if $m < k$ the next set $U_{m+1}$ will fulfill the first condition. 
	    	
	    	Now assume $U_m \cap \{1, 2, \cdots, M\} = \{2^{m-1}c, c = 1, 2, \cdots, \lfloor \frac{M}{2^{m-1}}\rfloor\}$, 
	    	in particular $T_{x_m} = 2^{m-1}$. 
	    	Since $G_{x_m}\in \{2^m - 2, 2^m - 1, 2^m\}$, this would give 
	    	$S_{x_m} \cap \{1, 2, \cdots, M\} = \{2^{m-1} + 2^{m}c, c = 0, 2, \cdots, \lfloor \frac{M - 2^{m-1}}{2^{m}}\rfloor\}$. 
	    	But the next element of $U_m$ after $M$ is $M + 2^{m-1}$, showing that 
	    	\[
	    	M + 2^{m-1} - 2^m \lfloor \frac{M - 2^{m-1}}{2^{m}}\rfloor\le \max G_{x_m}\le 2^m
	    	\]
	    	So the remainder of $M - 2^{m-1}$ when divided by $2^m$ must be less than $2^{m-1}$, 
	    	which, combined with previous evidence, gives $M\equiv 2^m-1\pmod{2^m}$. 
	    	This establishes that $U_m\cap \{1, 2, \cdots, M\}$ has odd number of elements (so $2^m \mid M + 1$) 
	    	and using $G_{x_m}\in \{2^m - 2, 2^m - 1, 2^m\}$, we iteratively have 
	    	$U_{m+1} = \{Mu + 2^{m}c, c = 1, 2, \cdots, \lfloor \frac{M}{2^{m}}\rfloor\}$ whenever $m < k$. 
	    \end{proof}
        Now to finish, we note that when $M = k$ we have $U_k = S_k$ so by the previous lemma, $U_k\cap \{1, 2, \cdots, M\}$ must have one element. 
        Since $2^k\mid M+1$ this is only possible when $2^k = M + 1$. 
        With this, all the $S_{x_i}$'s are uniquely determined and we may show that they map exactly to the number $x_m = \frac{2^{k-m}}{2^k-1}$. 
		
	\end{enumerate}
    
    \section*{Combinatorics}
    \begin{enumerate}
    	\item [C1.]
    	A line in the plane is called $sunny$ if it is not parallel to any of the $x$–axis, the $y$–axis, or the line $x+y=0$.
    	
    	Let $n\ge3$ be a given integer. Determine all nonnegative integers $k$ such that there exist $n$ distinct lines in the plane satisfying both of the following:
    	\begin{itemize}
    		\item for all positive integers $a$ and $b$ with $a+b\le n+1$, the point $(a,b)$ lies on at least one of the lines; and
    		\item exactly $k$ of the $n$ lines are sunny.
    	\end{itemize}
    	
    	\textbf{Answer.} $k=0, 1, 3$. 
    	
    	\textbf{Solution.} Let $T_n$ be the set of allowed $k$'s. The main claim is the following. 
    	
    	\textit{$T_n$ is independent of $n$ for all $n\ge 3$. }
    	To see why, we easily have $T_n\subseteq T_{n+1}$: just add a line (non-sunny) on the line $x+y=n$. 
    	To show $T_{n+1}\subseteq T_n$: consider the outermost triangle: union of points $x+y=n, x=0$ and $y=0$. 
    	We claim that one of the lines must coincide with one the three sides, which reduces to the case of $n$ lines. 
    	Suppose otherwise, then: 
    	\begin{itemize}
    		\item Each line not under this category has at most two intersects with the three sides in total (e.g. Menelaus') theorem.
    		covers at most 2 lines in the outer triangle. Jointly, this gives at most $2(n+1)$ intersections. 
    		
    		\item On the other hand, there are exactly $3n$ points on the outer triangle, and $3n > 2(n+1)$. 
    	\end{itemize}
    	
    	Thus it suffices to consider $n = 3$. We do three cases here. 
    	
    	\textit{$k=0$ or 1 construction.} $x=0, x=1$, and any line through $(2, 0)$, which we could choose to be sunny or not.
    	
    	\textit{$k=3$ construction.} $\{(0, 2), (1, 0)\}, \{(2, 0), (0, 1)\}, \{(0, 0), (1, 1)\}$. 
    	
    	\textit{Why $k=2$ fails.} We split by cases: if one of the three lines are $x=0, y=0$ or $x+y=2$ then one of other two lines (covering the smaller triangle) must be non-sunny (to cover 2 of out the 3 points). 
    	Otherwise, no line can cover more than 2 points (hence each line covers exactly two points). 
    	Can verify that the three lines of the $k=3$ case are the only sunny lines through exactly two points, 
    	hence choosing 2 of them means we also need the third. 
    	
    	
    	\item [C2.]
    	There is a row of $n$ paddocks, labelled from left to right with the integers $1$ to $n$, where $n\geq 3$. Skippy the Kangaroo grazes in the paddock according to the following rule:
    	
    	If Skippy is grazing in paddock $k$, then she makes a sequence of $k$ hops from a paddock to an adjacent paddock. The first of the $k$ hops is always to the right, unless Skippy is in paddock $n$, in which case it is to the left. Each of the following $k-1$ hops is in the same direction as the previous hop, unless Skippy is in paddock $1$ or $n$. Skippy grazes in the paddock that she is in after the $k^{\text{th}}$ hop.
    	
    	For example, if $n=8$ and Skippy is grazing in paddock $3$, the next $4$ paddocks she will graze in are $6$, $4$, $8$, $2$, in this order. Skippy continues grazing in this way indefinitely. A paddock is overgrazed if Skippy eventually grazes in it, regardless of which paddock she starts in. Determine all $n\geq 3$ for which there exists exactly one overgrazed paddock.
    	
    	\textbf{Answer.} $n=3\cdot 2^k$. 
    	
    	\textbf{Solution.} Denote $f(k)$ as the next grazing grid of Skippy when she is grazing on $k$. If $k$ is overgrazed, so is $f(k), f(f(k)), \cdots$. Therefore if $k$ is unique overgrazer, $k=f(k)$. 
    	We now note the following rule:
    	\[
    	f(k) = 
    	\begin{cases}
    		2k & k \le \frac{n}{2}\\
    		2(n-k) & \frac{n}{2} < k \le n - 1\\
    		2 & k = 2\\
    	\end{cases}
    	\]
    	We start with these observations: 
    	\begin{itemize}
    		\item If $k = f(k)$, then $k = \frac{2n}{3}$, or $3\mid n$. 
    		
    		\item If $k\neq n$, and $\nu_2(k)\le \nu_2(n)$, then $\nu_2(f(k)) > \nu_2(k)$. 
    		
    		\item If $k\neq n$, $p > 1$ odd, $p\mid n$, then $p\mid k$ if and only if $p\mid f(k)$. 
    	\end{itemize}
        In particular, the first bullet point allows us to disregard powers of 2. 
        The next we will eliminate is where $n$ is divisible by $pq$ where $p, q > 1$ are odd. 
        The last bullet point suggests that Skippy will only cycle among grids not divisible by $p$ when she starts at 1. 
        However, if she starts at $p$, the subsequent grids she stays at cannot be divisible by $pq$, 
        and therefore cannot land at $n$ at all. 
        On the other hand, this also means all the grids it stays are divisible by $p$. 
        Thus the two paths has two terminal cycles that are disjoint, showing that there is no overgrazed pollock for tis $n$. 
        
        We are left with the case $n=p\cdot 2^k$, where $p$ is prime. 
        Given the requirement $3\mid n$, we have $n = 3\cdot 2^k$ for $k\ge 0$. 
        We now claim $2^{k+1}$ is the only overgrazed pollock: 
        in fact, all paths would terminate here.
        When $m=2^{k+1}, f(m)=m$. 
        When $m=n$, $f(m)=2$ and eventually leads to $2^{k+1}$. 
        In all other cases $\nu_2(f(m))\ge \nu_2(m) + 1$, so $\nu_2$ increases as Skippy goes 
        (possibly hitting $m=n$ but then goes back to the cycle to $2^{k+1}$. )
    	
    	\item [C3.]
    	There are $2025$ white balls and $2025$ black balls arranged in a row. A balanced segment is a contiguous nonempty segment of balls that contains the same number of white balls as black balls. A balanced removal is the operation of selecting a balanced segment $S$, removing all the balls in $S$, and shifting left the balls that were to the right of $S$ so that the remaining balls form a contiguous row.
    	
    	Determine the smallest positive integer $N$ satisfying the following property: for every configuration of balls and for every integer $k$ satisfying $0\le k\le2025$, there exists a sequence of at most $N$ balanced removals after which there are exactly $2k$ remaining balls.
    	
    	\textbf{Answer.} $N = 2$. 
    	
    	\textbf{Solution.} We generalize to $n$ white and $n$ black balls, $n\ge 4$. 
    	To see why $N = 1$ fails, 
    	consider any configuration when first 2 and last 2 balls are all white; then $k = 1$ fails here. 
    	
    	To see why $N = 2$ always works, we consider the notion of `prefix' as the subset of first $m$ balls for any $m$; call this prefix strict if $m < 2k$; denote suffix similarly (as consecutive $m$ balls). 
    	We start with the following observation:
    	
    	\textit{N = 1 works when any strict prefix has more white than black balls (or vice versa).} W.l.o.g. we only consider the first case, and the task of removing $2(n-k)$ consecutive balls. 
    	Let $m_g$ be the difference between white and black balls among balls $g+1, \cdots, g+ 2(n-k)$ for $g=1, 2, \cdots, 2k$. 
    	Then $m_{2k} < 0 < m_0$, $m_i$'s are all even, and morever $|m_g-m_{g+1}|\in \{0, 2\}$. 
    	Hence the conclusion follows by discrete continuous value theorem. 
    	
    	To finish, consider the smallest prefix $P$ of at least $2k$ balls and is balanced; the first move would be remove the suffix of the whole sequence, with $P$ left. 
    	Let $Q$ be the smallest suffix of $P$ that is balanced, and let $m=|P|-2k > |Q|$ be the number of balls we need to remove. By minimality, the condition of our previous claim applies, 
    	hence we may remove $m$ balls from this suffix $Q$ in one step. 
    	
    	\item [C4.]
    	Bluey is placing lollies on a $1000\times1000$ grid. Initially, every cell of the grid is empty. At each step, Bluey chooses a row or column in which the total number of lollies is a multiple of $3$, chooses an empty cell in this row or column, and places a lolly in it. If there is no such row or column, Bluey stops.
    	
    	Determine the maximum number of lollies that Bluey can place.
    	
    	\textbf{Answer.} 3995 (or $4n-5$ for all $n\ge 5$). 
    	
    	\textbf{Solution.}
    	\textit{Bound:} let the `score' of each row and column be 0, 1, 2 when the number of lollies is congruent to 0, 2, 1, respectively. 
    	Denote total score as total row score and total column score across all $n$ rows and columns 
    	(i.e. range 0 to $4n$). Consider each of the lolly operation (modulo 3) : 
    	\[
    	(0, 0)\to (1, 1), (0, 1) \to (1, 2), (0, 2) \to (1, 0)
    	\]
    	Call this type 0, 1, 2. 
    	Corresponding score update is $+4, +1, +1$, respectively, 
    	and the first update must be type 0 (hence +4 points). 
    	Now the last update: 
    	\begin{itemize}
    		\item If it's type 0, max total number of steps is $4n-8+2=4n-6$; 
    		
    		\item If type 1 or 2, then there is row (and column) where the score is $\le 2$. 
    		The total row score and total column score must align modulo 3 (hence $\le 2n-1$ each); total score $\le 4n-2$ so including the first 4-point update total number of steps is $\le 4n-5$. 
    	\end{itemize}
    	
    	\textit{Construction:} first do 
    	\[
    	(1, 1), (i, i+1), (i+1, i), i=1, 2, \cdots, n-2
    	\]
    	Now consider permutation $\pi$ such that $\pi(1)=n$, $\pi(n)=n-1$
    	and $|\pi(i)-\pi(i+1)|\neq 1$ for all $i\ge 1$. 
    	E.g. for $n$ even we may do 
    	\[
    	1, 3, \cdots, n - 3, n, n - 2, \cdots, 2, n - 1
    	\]
    	for $n$ odd: 
    	\[
    	1, 3, \cdots, n, 2, 4, \cdots, n-1
    	\]
    	Then row/column $\pi(1)$ has 0 lolly, $\pi(n)$ has 1, and $\pi(i)$ has 2 for $i=2, \cdots, n-1$. 
    	Then do 
    	\[
    	(\pi(i), \pi(i+1)), (\pi(i+1), \pi(i)), i=1, 2, \cdots, n-1
    	\]
    	which gives total of $4n-5$ operations. 
    	
    	\item [C5.]
    	Let $a_1,a_2,\dots,a_n$ be a permutation of $1,2,\dots,n$. A pair $(i,j)$ with $i<j$ is called an inversion if $a_i>a_j$. An inversion is called odd if $i-j$ is odd.
    	
    	\begin{enumerate}
    		\item[(a)] Prove that the total number of inversions is at most $5$ times the number of odd inversions.
    		
    		\item [(b)]There is a permutation such that the total number of inversions is more than $4.99$ times the number of odd inversions.
    	\end{enumerate}
        
        \textbf{Solution.} 
        
        (a)
        Define `even inversion' similarly. 
        Note that for any inversion $(i, j)$ with $i < j$, 
        any $k$ with $i < k < j$, one (or both) of $(i, k)$ and $(k, j)$ must also be an inversion. 
        
        We construct a bipartite graph with one side even inversion, 
        and the other side odd inversion, in the following manner:
        for every even inversion $(i, j)$ with $j - i = 2d > 0$, 
        if $d$ is odd, 
        we connect an edge of weight 1 to either $(i, i + d)$ or $(i + d, j)$, whichever is an inversion; 
        if $d$ is even, 
        we consider one of $(i, i + d - 1), (i + d - 1, j)$ that is an inversion,
        and also one of $(i, i + d + 1), (i + d + 1, j)$ that is an inversion, 
        and connect $(i, j)$ to each of the chosen inversions with weight $\frac 12$. 
        Now note that each even inversion is incident edges of total weight 1, thus the total weight of edges is $A - B$. 
        Meanwhile for each odd inversion $(i, j)$ with $j = i + d$, the possible neighbours are: 
        \[
        (i, j + d), (i - d, j), (i, j + d + 1), (i - d - 1, j), 
        (i, j + d - 1), (i - d + 1, j)
        \]
        The first two (if chosen) have weights 1 each and the remaining four have weight $\frac 12$ each, 
        so each of odd inversions are incident to edges of total weight at most $2 + 4\cdot 12 = 4$. 
        It follows that we have $A - B \le 4B$ or $A\le 5B$. 
        
        (b) Let $k, m > 1$ and $n = 2km$. 
        We consider $m$ blocks, each $2k$ elements. 
        Then organize the $n$ numbers as follows: 
        \begin{itemize}
        	\item We consider $2m$ classes of $k$ numbers each, 
        	if $i < j$ then elements in class $i$ is strictly smaller than elements in class $j$ 
        	(w.l.o.g. can treat the $k$ elements of class $j$ as $jk, jk + 1, \cdots, (j+1)k-1$). 
        	
        	\item Each block gets two classes; the only inversion at block level is that for two adjacent blocks $i$ and $i + 1$, 
        	bigger class at block $i$ is bigger than smaller class at block $i + 1$. 
        	
        	\item Within each block $i$ of classes $a, b$ ($a < b$), 
        	the numbers are presented in alternating and decreasing order; 
        	if $i$ is odd then odd indices within this block take class $a$ and even indices take class $b$. 
        	Else the opposite. 
        \end{itemize}
        E.g. when $k = m = 4$, the classes will take the form 
        \[
        (0, 2), (1, 4), (3, 6), (5, 7)
        \]
        The first block will have number (denoted in (class, number) for brevity so $(a, b)$ means $ak + b$)
        \[
        (0, 3), (2, 3), (0, 2), (2, 2), (0, 1), (2, 1), (0, 0), (2, 0)
        \]
        and the second block 
        \[
        (4, 3), (1, 3), (4, 2), (1, 2), (4, 1), (1, 1), (4, 0), (1, 0)
        \]
        
        We now claim that when $m$ is even, $B = \frac{k^2m}{2}$, 
        and $A - B = k(k-1)m + k^2(m-1)$. 
        
        Now, inversion involving pair of different blocks could only happen for adjacent blocks: bigger class at block $i$ vs smaller class at block $i + 1$ ($i=1, 2, \cdots, m - 1$); 
        by construction, such inversions must be even and contributes to $k^2$ each, thus contributing $k^2(m-1)$ to $A-B$. 
        
        Within the same block, the two sets of $k$ numbers of the same class are of even inversion, thereby contributing to $2\binom{k}{2}=k(k-1)$ inversions. 
        Taken together, this contributes another $k(k-1)m$ to $A-B$. 
        Finally, for the same block $i$ but different classes, 
        around half are odd inversions: 
        $\binom{k}{2}$ for $i$ odd but $\binom{k+1}{2}$ for $i$ even. 
        This gives a total of $\frac{k^2m}{2}$ odd inversions. 
        Therefore we have 
        \[
        \frac{A}{B}
        = 1 + \frac{k(k-1)m + k^2(m-1)}{k^2m/2}
        = 1 + \frac{2(k-1)}{k} + \frac{2(m-1)}{m}
        = 5 - \frac{2}{k} - \frac{2}{m}
        \]
        Taking $k = m = 1000$ solves the problem. 
    	
    	\item [C7.]
    	Let $P$ be a regular $100-$gon. An Aussie triangulation is formed by dividing $P$ into a set $T$ of non-overlapping triangles such that:
    	\begin{itemize}
    	\item there are exactly $2025$ different points, including the $100$ vertices of $P$ that are a vertex of at least one triangle in $T$, and
    	
    	\item no three of these $2025$ points are collinear.
    	\end{itemize}
    	A quokkalateral is a convex quadrilateral that consists of two triangles in $T$ sharing a common side. Two quokkalaterals may share a triangle in $T$. Determine the minimum number of quokkalateral among all Aussie triangulation.
    	
    	\textbf{Answer.}  1011. 
    	
    	\textbf{Solution.} 
    	
    	Replacing 2025 and 100 with $n$ and $m$ respectively, with $n > 2(m - 1)$ and $n$ odd, we show that in the configuration with $n$ points where the convex hull has exactly $m$ points has a minimum of ? quokkalaterals. 
    	Call the $n - m$ points inside the convex hull the 	internal' points, and `internal' lines as all lines forming the triangle other than the $m$ sides of the convex hull. 
    	Denote the \textit{candidate quadrilaterals} as quadrilaterals that has two triangles sharing a common side; 
    	a quokkalateral is a candidate quadrilateral. 
    	
    	To begin, consider following setup: 
    	
    	\begin{itemize}
    		\item 
    		\textit{Number of triangles}. 
    		Each internal point contributes total angle of $2\pi$ to the triangles. Together with the $(m-2)\pi$ contributed by the convex hull, the number of triangles is $(m - 2) + 2(n - m) = 2n - m - 2$. 
    		
    		\item \textit{Number of internal lines.} Each convex hull side and internal line corresponds to 1 and 2 triangles, 
    		respectively. Consequently, the number of internal lines are 
    		$\frac{3(2n-m-2) - m}{2} = 3n - 2m - 3$. 
    		
    		\item \textit{Line-to-quadrilateral bijection.} 
    		Each candidate quadrilateral corresponds to exactly one internal line (i.e. there is a bijection between the two). 
    		
    		\item \textit{Counting and minimizing quokkalateral.} 
    		Consider any candidate quadrilateral; we need to check the internal lines in $T$ correspond to convex angles to determine if it is quokkalateral. 
    		In other words, minimizing quokkalaterals (among all the allowable configurations) is equivalent to maximizing internal lines in $T$ that has exactly one endpoint where the triangles adjacent to it forms a reflex angle. 
    	\end{itemize}
        In essence, we can also count now many reflex angles are contributed by each point in $T$; such points have to be internal (not the $m$ points on the convex hull). 
        
        Before proceeding, we define a \textit{supertraingle} as a collection of 3 triangles $OAB, OAC, OBC$ such that $O$ is inside triangle $ABC$; call $O$ its center for this purpose. The main claim is the following. 
        
        \begin{lemma}
        	An internal point contributes to exactly 3 reflex angles if and only if it is the center of a supertriangle. 
        	Moreover, an internal point contributing to exactly two reflex angles must correspond to a vertex of a triangle in $T$ that cannot be part of a supertriangle; 
        	each such internal point must correspond to different triangle in $T$. 
        \end{lemma}
        
        \begin{proof}
        	For the first part, the `if' part is clear. To solve for `only if' part, if a point $O$ corresponds to at least two reflex angles, 
        	these two reflex angles must overlap, 
        	hence share a triangle in $T$. The only possibility when each of the 3 pairs of reflex angles overlap is when this point $O$ is incident to exactly three internal lines, which then forms a supertriangle. 
        	
        	For the second part, consider the angle shared between two reflex angles incident to $O$. This angle corresponds to  triangle in $T$ where $O$ is a vertex. 
        	Now, $O$ cannot be the center of a supertriangle. 
        	Each the two regions in $O$ adjacent to this angle add up to a reflex angle with this region, hence $O$ cannot be part of a supertriangle either. 
        	Finally, suppose this triangle also corresponds to anther point $R$ with two reflex angles. 
        	This point $R$ must also be a vertex of this triangle and the line $OR$ has both endpoints ($O$ and $R$) corresponding to a reflex angle. 
        	This is impossible as no quadrilaterals can have two reflex angles. 
        \end{proof}
        
        
        
        Now, let $k_3$ be the number of supertriangles, 
        $k_2$ number of points contributing to two reflex angles, 
        and $k_1$ the number of remaining points. Then $k_1+k_2+k_3=n - m$ and: 
        \[
        k_3\le \min \{\frac{2n-m-2}{3}, n - m\}
        \qquad 
        k_2\le \min \{2n - m - 2 - 3k_3, n - m - k_3\}
        \]
        where $2n-m-2$ denotes the total number of triangles, 
        and $2n - m - 2 - 3k_3$ denotes the non super triangles. 
        Then the number of nonconvexity is upper bounded by 
        \[
        3k_3+2k_2+k_1 = 2k_3 + k_2 + (n - m)
        \]
        We note that the function $2x + \min \{2n - m - 2 - 3x, n - m - x\}$ is piecewise linear in $x$ that attains maximum value when $2n - m - 2 - 3x = n - m - x$, i.e. when $x = \frac{n - 2}{2}$. 
        Since $n \ge 2m - 2$, $\frac{n-2}{2}\le \min \{\frac{2n-m-2}{3}, n - m\}$. 
        However, our $n$ is odd, hence we need to consider the adjacent values $\frac{n-3}{2}$ and $\frac{n-1}{2}$; 
        both give the value $\frac{3(n-1)}{2}-m$. 
        Thus the number of quokkalateral is at least 
        \[
        3n - 2m - 3 - (\frac{3(n-1)}{2} - m) - (n - m)
        =\frac{n - 3}{2}
        \]
        
        Construction: (another time, lol)
    	
    \end{enumerate}
    
    \section*{Geometry}
    \begin{enumerate}
    	\item [G1.]
    	Let $ABC$ be a triangle such that $\angle A$ is obtuse. Let $D$ and $E$ be points on sides $AB$ and $AC$, respectively, such that $BDEC$ is cyclic. Suppose the circumcircle of triangle $ADE$ intersects side $BC$ at two points $X$ and $Y$, where $B$, $X$, $Y$ and $C$ lie in that order. Let the circumcircles of triangles $BDX$ and $CEY$ be $\Omega$ and $\Gamma$, respectively. Let $P$ be a point such that $PD$ is tangent to $\Omega$ and $PE$ is tangent to $\Gamma$. Let $Q$ be a point such that $QB$ is tangent to $\Omega$ and $QC$ is tangent to $\Gamma$. Prove that points $A$, $P$ and $Q$ are collinear.
    	
    	\textbf{Solution.} We show that $A, P, Q$ all lie on radical axis of $\Omega$ and $\Gamma$. 
    	
    	$A$: power to $\Omega$ and $\Gamma$ are $AD\cdot AB$ and $AE\cdot AC$, respectively. 
    	The two quantities are equal because $B, D, C, E$ cyclic. 
    	
    	$Q$: we claim that $BQ=QC$. Note that if $O$ is circumcenter of $ADE$ 
    	then $\angle DAO = 90^{\circ} - \angle AED = 90^{\circ} - \angle ABC$ so $AO\perp BC$. 
    	Thus $AX=AY$ and 
    	$\angle CBQ = \angle BDX=\angle CEY=\angle QCB$. 
    	
    	$P$: similarly, we show $PD=PE$. 
    	One sees that $\angle AED=\angle ABC=\angle XDP$ and $\angle ADE=\angle ACB=\angle PEY$, 
    	and $\angle BDX=\angle CEY$ as previously claimed. 
    	
    	\item[G2.] Let $ABC$ be an acute-angled triangle. Let $I_B$ and $I_C$ be the excenters of triangle $ABC$ opposite $B$ and $C$, respectively. Let $E$ be a point on line $BA$ such that $\angle EI_BB=90^\circ$. Let $F$ be a point on line $CA$ such that $\angle CI_CF=90^\circ$. Let the circle with center $I_B$ and radius $I_BE$ intersect the segment $I_BB$ at $X$. Let the circle with center $I_C$ and radius $I_CF$ intersect the segment $I_CC$ at $Y$. Prove that $BY$ is perpendicular to $CX$.
    	
    	\textbf{Solution.} 
    	The main claim is that triangles $I_CBY$ and $I_BXC$ are similar, 
    	and that both angles $\angle BI_CY=\angle CI_BX=\frac{\angle BAC}{2}$. 
    	Assuming the claim, 
    	we then have 
    	\[
    	\angle YBC +\angle XCB
    	=\angle I_CBC+\angle I_BCB-\angle I_CBY-\angle I_BCX
    	\]\[
    	= (90^{\circ}+\frac{\angle ABC}{2}) + (90^{\circ}+\frac{\angle ACB}{2}) - (180^{\circ}-\frac{\angle BAC}{2})
    	=90^{\circ}
    	\]
    	where the last equality follows from that the sum of 3 angles in a triangle is $180^{\circ}$. 
    	
    	To prove the claim, the angle claim $\angle BI_CY=\angle CI_BX=\frac{\angle BAC}{2}$ follows directly from angle chasing. 
    	In particular, $\angle BI_CY=180^{\circ} - \frac{\angle ACB}{2} - (90^{\circ} + \frac{\angle ABC}{2})
    	=\frac{\angle BAC}{2}$. 
    	Thus, proving similarity of the triangle now reduces to proving that $\frac{I_CB}{I_CY}=\frac{I_BX}{I_BC}$. 
    	Now $I_CY = I_CF = I_CC\tan \frac{\angle ACB}{2}$; 
    	by sine rule we have $I_CB=I_CC\frac{\sin\frac{\angle BCA}{2}}{\cos \angle \frac{ABC}{2}}$. 
    	Therefore $\frac{I_CB}{I_CY}=\frac{\cos \angle \frac{ACB}{2}}{\cos \angle \frac{ABC}{2}}$. 
    	The other equality $\frac{I_BX}{I_BC}=\frac{\cos \angle \frac{ACB}{2}}{\cos \angle \frac{ABC}{2}}$ can be established similarly. 
    	
    	\item [G3.] Let $ABC$ be an acute triangle with $AB < AC$. Let $\Gamma$ be the circumcircle of $ABC$. The perpendicular from $B$ to $AC$ intersects $\Gamma$ again at $D\neq B$. Let $Q$ be the point on $\Gamma$, different from $B$, such that $AB=AQ$. Lines $BQ$ and $AC$ intersect at $R$. Line $DR$ intersects $\Gamma$ again at $S\neq D$. Segment $BC$ intersects the circumcircle of triangle $QRS$ at $T$. Lines $TQ$ and $RS$ intersect at $X$. Prove that $CX$ is perpendicular to $AB$.
    	
    	\textbf{Solution.} 
    	We first show that $QT\perp AC$. 
    	Let $T'$ on $BC$ such that $QT'\perp AC$. 
    	Since $AB=AQ$, $\angle ACB=\angle ACQ$, so $QT'$ is perpendicular bisector of $AC$. 
    	This then gives the following sequence of equalities: 
    	\[
    	\angle RT'Q=\angle RQT'=\angle DBQ=\angle DSQ
    	\]
    	The second equality due to $BD\parallel QT'$ and last equality due to $D, B, S, Q$ concylic on $\Gamma$. 
    	Therefore, $T'$ is on the circle $QRS$ and on segment $BC$. 
    	Given that the other intersection of circle $QRS$ and line $BC$ will be on its extension beyond $C$, 
    	we have $T=T'$. 
    	
    	Next, let $AC$ intersect $BD$ at $U$, and $QT$ intersect $AC$ at $V$. 
    	Then given $BD\parallel QT$ and $D, R, X$ collinear, $\frac{BU}{UD}=\frac{VQ}{VX}$. 
    	This then gives $\frac{\tan\angle CAB}{\tan\angle CAD}=\frac{\tan\angle ACQ}{\tan\angle ACX}$, 
    	and given that $\angle CAD=90^{\circ}-\angle ACB=90^{\circ}-\angle ACQ$, 
    	the other two angles $\angle CAB$ ad $\angle ACX$ also add up to $90^{\circ}$. 
    	This give $CX\perp AB$, as required.  
    	
    	\item [G4.] (IMO 2) 
    	Let $\Omega$ and $\Gamma$ be circles with centres $M$ and $N$, respectively, such that the radius of $\Omega$ is less than the radius of $\Gamma$. Suppose $\Omega$ and $\Gamma$ intersect at two distinct points $A$ and $B$. Line $MN$ intersects $\Omega$ at $C$ and $\Gamma$ at $D$, so that $C, M, N, D$ lie on $MN$ in that order. Let $P$ be the circumcentre of triangle $ACD$. Line $AP$ meets $\Omega$ again at $E\neq A$ and meets $\Gamma$ again at $F\neq A$. Let $H$ be the orthocentre of triangle $PMN$.
    	
    	Prove that the line through $H$ parallel to $AP$ is tangent to the circumcircle of triangle $BEF$.
    	
    	\textbf{Solution.} 
    	Claim: $P$ is $A$-excenter of triangle $AMN$. 
    	
    	Proof: $CM=MA$ and $PC=PA$ so $MP$ is perpendicular bisector of $AC$ (so $MP$ bisects $\angle CMA$), 
    	same for $NP$ and $\angle DNA$. 
    	
    	Introduce $E_1$ midpoint of $AE$, $F_1$ midpoint of $AF$, and $H$ midpoint of $AB$. Let circle $HE_1F_1$ intersect $CD$ again at $G$, 
    	and $I$ incneter of triangle $AMN$. 
    	Then: 
    	\begin{itemize}
    		\item $H$ is foot of perpendicular from $A$ to $CD$; 
    		\item $CE_1\perp AP$ and $DF_1\perp AP$; 
    		\item $HE_1AC$ and $HF_1AD$ cyclic; 
    		\item $\angle E_1HG=\angle CAP=\angle DAP=\angle F_1HG$, so $E_1G=F_1G$ 
    		\item $\frac{E_1G}{CG} = \frac{\sin\angle E_1CD}{\sin\angle CE_1G}$ and 
    		$\frac{F_1G}{DG} = \frac{\sin\angle F_1DC}{\sin\angle DF_1G}$. 
    		Note $\angle CE_1G+\angle DF_1G=180^{\circ}$ and $\angle E_1CD = \angle F_1DC$, and $E_1G=F_1G$ so $G$ is midpoint of $CD$. 
    		\item Finally, $IMNH$ is parallelogram. 
    	\end{itemize}
    	Thus to finish: 
    	\begin{itemize}
    		\item Tangent at $G$ to circle $HE_1F_1$ is parallel to $AP$, call this line $\ell_0$. 
    		
    		\item Since $G$ is also midpoint of $IH$, line $\ell_1$ through $H$ parallel to $AP$ is $2\times$ distance to $AP$ as compared to line $\ell_0$. 
    		
    		\item $BEF$ is $2\times$ homothety compared to $HE_1F_1$ with center $A$, so $\ell_1$ is tangent to circle $BEF$. 
    	\end{itemize}
        
        \item [G5.] Let $ABCD$ be a convex quadrilateral with $\angle ADC > 90^{\circ}$. Suppose that points $X$ and $Y$ lie on diagonal $AC$ such that
        $$\angle ADX = \angle YDC = 90^{\circ}, \quad \angle CAD = \angle CBX, \quad \text{and} \quad \angle DCA = \angle YBA.$$Let $O$ be the circumcentre of triangle $BXY$. Suppose point $R \neq O$ lies on the perpendicular bisector of $AC$ such that lines $OR$ and $AC$ are parallel.
        
        Prove that line $RO$ bisects angle $\angle DRB$.
        
        \textbf{Solution.} 
        Let the circumcircle of $DAC$ intersect $DY$ and $DX$ again at $P, Q$, respectively. 
        Then $APQC$ is a rectangle, 
        and in addition $\angle YBA=\angle DCA=\angle DPA$, 
        so $AYBP$ is cyclic and similarly $CXBQ$ is cyclic. 
        Consequently, $\angle YBP=\angle YAP=90^{\circ}$ and similarly $\angle XBQ=\angle XCQ=90^{\circ}$. 
        
        We now claim that $D, O, B$ are collinear. 
        In fact, let $E$ be such that $\angle EYB=\angle EXB=90^{\circ}$, 
        then $O$ is the midpoint of $EB$. 
        Meanwhile, $EY\parallel BP$ and $EX\parallel BQ$, so 
        $D, E, B$ are also collinear, and thus both $D$ and $O$ are on the line $EB$. 
        
        Denote $S$ as the point on line $DB$ such that $PS=SQ$; 
        this is obtained by sending $SP\parallel YO$ and $SQ\parallel OX$. 
        Observe that $S$ and $R$ are both on the perpendicular bisector of $AC$. 
        We now claim that $(D, B; O, S)$ is a harmonic bundle, 
        i.e. $\frac{OB}{OD}=\frac{SB}{SD}$. 
        Now, the homothety map $\Gamma$ centered at $D$ that brings $Y$ to $P$ has $\Gamma(X)=Q$, $\Gamma(E)=B$ and $\Gamma(O)=S$. 
        Therefore, $\frac{SB}{SD}=\frac{OE}{OD}$. 
        Given that $OE=OB$ as $O$ is midpoint of $EB$,  this establishes our claim. 
        
        Finally, since $OR\perp RS$ and $(D, B; O, S)$ harmonic, 
        it follows that $OR$ and $RS$ are internal and external bisectors of $\angle DRB$, respectively. 
        
    \end{enumerate}
    
    \section*{Number Theory}
    \begin{enumerate}
    	\item [N1.]
    	Let $n$ and $k$ be positive integers such that $n < k < 2n$. Suppose $a_1, a_2, \ldots, a_k$ are positive integers such that $2^{a_1} + 2^{a_2} + \cdots + 2^{a_k}$ is divisible by $2^n - 1$.
    	
    	Show that at least three of $a_1, a_2, \ldots, a_k$ have the same remainder when divided by $n$.
    	
    	\textbf{Solution.} Take all exponents modulo $n$. 
    	Let $a = \sum 2^{b_j}$ and $b = \sum 2^{c_j}$, 
    	$b_j$ are different indices occuring exactly twice and $c_j$ exactly once. 
    	Then $2^n-1 \mid 2a + b$, 
    	and $a, b$ have disjoint positions of 1s (binary form). 
    	$a, b$ cannot be 0 given $n < k < 2n$. 
    	
    	Case $b$ odd: $2a+b=2^n-1$ and $2a, b$ have complementary digits at each place: 
    	rightmost 0 of $b$ = 1 for $2a$, its right is 1 for both $a$ and $b$, contradiction. 
    	
    	Case $b$ even: $2a+b=2(2^n-1)$. $a$ and $\frac{b}{2}$ have complementary digits at each place: 
    	leftmost 1 of $\frac{b}{2}$: to the left is 0 for $\frac{b}{2}$ but 1 for both $a$ and $b$, contradiction. 
    	
    	\item[N2.]
    	In an $n \times n$ board, for each $1 \leq i \leq n$ and $1 \leq j \leq n$, the cell in the $i$th row from
    	the bottom and $j$th column from the left contains $\gcd(i,j)$ leaves. A koala travels from the bottom left
    	corner to the top right corner. At each step, the koala moves up or to the right by a single cell. The
    	koala eats the leaves in each cell it visits, including the bottom left and top right cells.
    	
    	Determine the maximum number of leaves that the koala can eat.
    	
    	\textbf{Answer.} $\frac{n(n+1)}{2} + (n - 1)$. 
    	
    	\textbf{Solution.} 
    	\textit{Construction.} 
    	This path is in zigzag manner (i.e. $(k-1, k)$ on odd-numbered step and $(k, k)$ on even-numbered step). 
    	
    	\textit{Bound.} 
    	Let the initial step index be 2 and ending step inded be $2n$. 
    	start at step 2 (w.l.o.g.) and end at step $2n$; 
    	the sum of coordinates at step indexed $i$ is $i$. 
    	We show that the total number of leaves at steps $2k-1$ and $2k$ cannot exceed $k + 1$. 
    	To proceed, one of the grids must have strictly more than $\frac{k}{2}$ leaves. 
    	This means we are down to one of the following 3 cases. 
    	
    	\textit{Step $2k$ contains $k$ leaves}. Then the coordinate at step $2k$ is $(k, k)$, 
    	at step $2k-1$ is either $(k-1, k)$ or $(k, k-1)$ and can only contain 1 leaf. 
    	
    	\textit{Step $2k$ contains $\frac{2k}{3}$ leaves}.By symmetry, we may assume that the coordinate is $(a, b)=(\frac{2k}{3}, \frac{4k}{3})$. 
    	Crucially, $b=2a$. Then the coordinate at step $2k-1$ is either $(a-1, 2a)$ or $(a, 2a+1)$, 
    	which has $\gcd\le 2$. 
    	Hence the total number of leaves is at most $2 + \frac{2k}{3}\le k + 1$, since $k\ge 3$. 
    	
    	\textit{Step $2k - 1$ contains $\frac{2k - 1}{3}$ leaves}. We can similarly deduce that the coordinates must be 
    	$(a', b')$ with $b'=2a'$, 
    	and that the coordinates at step $2k$ is either $(a', 2a'+1)$ or $(a'+1, 2a')$. 
        Either way, the gcd is at most 2 and total number of leaves is $\le 2 + \frac{2k-1}{3}\le k + 1$ for $k\ge 2$. 
    	
    	\item [N3.] A proper divisor of a positive integer $N$ is a positive divisor of $N$ other than $N$ itself.
    	
    	The infinite sequence $a_1, a_2, \cdots$ consists of positive integers, each of which has at least three proper divisors. For each $n\geq 1$, the integer $a_{n+1}$ is the sum of the three largest proper divisors of $a_n$.
    	
    	Determine all possible values of $a_1$.
    	
    	\textbf{Answer.} $a_1 = 12^{k}\cdot 6c$, $c$ relatively prime to 2 and 5. 
    	
    	\textbf{Solution.} The main insight is that at no point can $a_n$ be monotonically strictly decreasing. 
    	This gives the following two claims. 
    	
    	\textit{$a_n$ must be even}. Otherwise, $a_{n+1}\le a_n(\frac{1}{3}+\frac{1}{5}+\frac{1}{7}) < a_n$ and $a_{n+1}$ is odd, so $a_m$ is monotonically decreasing for all $m\ge n$. 
    	
    	\textit{$a_{n}$ must be divisible by 3.} Suppose otherwise. 
    	Then $a_{n+1}=a_n(\frac{1}{a}+\frac{1}{b}+\frac{1}{c})$ where $a<b<c$ are the smallest divisors less than $n$, 
    	and greater than 1. 
    	Note that $a\ge 2, b\ge 4, c\ge 5$, so $a_{n+1}\le \frac{19}{20}a_n < a_n$. 
    	Moreover, we assumed $a_n$ is even, so $a=2$. 
    	If $b = 4$, $3\mid a_n(\frac{1}{a}+\frac{1}{b})$ so $3\nmid a_{n+1}$. 
    	If $b > 4$, $4\nmid a_n$, $b$ the smallest odd divisor of $a_n$ (i.e. odd prime). 
    	Then $c = 2b$. Now can show $a_{n+1}=a_n\cdot \frac{b+3}{2b}$ which is not divisible by 3 (as $3\nmid b$), so this cycle continues and $a_m$ is monotonically decreasing for all $m\ge n$. 
    	
    	Now with these set up, we now have the three smallest divisors greater than 1 as $2, 3, c$ where $c\in \{4, 5, 6\}$. 
    	Note, however, the case $c=5$ can't occur as this means $a_n$ is divisible by 2 but not 4, and so $a_{n+1}=\frac{31}{30}a_n$ is odd, which has been shown to be impossible. 
    	Call a number in this form type-$c$; $a_n$ of type 4, 6 gives $a_{n+1}=\frac{13}{12}a_n, a_n$ respetively. 
    	Since 13 and 31 are relatively prime to 4, 5, 6, so $a_{n+1}$ can never become a smaller type than $a_n$. 
    	It follows that $a_1$ has form $12^{k}A$ where $6\mid A$ but $4, 5\nmid A$. 
    	The terminating number in the sequence is then $13^kA$. 
    	
    	
    	\item [N4.]
    	Find all functions $f:\mathbb{N}\to \mathbb{Z}$ such that for every positive integers $n$ and $d$ such that $d\mid n$, there exits a positive integer $e\mid n$ such that $n\mid d+f(e)$.
    	
    	\textbf{Answer.} $f(x)=-x$. This fulfills the condition by taking $e = d$. 
    	
    	\textbf{Solution.} The equivalent statement of the problem condition is 
    	for each $n$, $\{f(d): d\mid n\} = \{-d: d\mid n\}\pmod{n}$. Denote this condition as (*). 
    	In particularly, (*) implies injectivity since $f(d_1)=f(d_2)$ implies $d_1=d_2$ when taking $n = kd_1d_2$ for $k=(|f(d_1)|+|f(d_2)|+1)$. 
    	The main claim is the following. 
    	
    	\begin{lemma}
    		For each $d > 0$, $f(d) < 0$ and $f(d)\mid m$. 
    	\end{lemma}
        
        \begin{proof}
    	For each $d, k\ge 1$ we choose $e_k$ as the number where $dk\mid f(d) + e_k$ and $e_k\mid dk$. 
    	Now choose $k$ where $dk > 2|f(d)|$. Note that either $\frac{1}{dk/e}=1$ or $\le \frac 12$, 
    	the first case implies $f(d)=0$; the second case implies $\frac{f(d)}{dk} = -\frac{1}{k'}<0$ for some $k'$ depending on $k$. 
    	Thus $f(d)\mid dk$ for all $k$  sufficiently large, 
    	so $f(d)\mid d$, 
    	
    	We're now left with showing $f(d)$ cannot be 0. 
    	By injectivity, there is at most one such $d_0$ with $f(d_0)=0$. 
    	If this $d_0$ exists, then this sequence of logic holds: 
    	\begin{itemize}
    		\item By injectivity, and $f(d)\mid n$ for $d\mid n$ but $d_0\nmid n$, we have $f(n)=-n$ for $n$ with $d_0\nmid n$. 
    		
    		\item By injectivity again we have $d_0\mid f(n)$ whenever $d_0\mid n$. 
    		
    		\item 
    		Consider $n = 2d_0, 3d_0$ simultaneously, however, $f(n)=-kd_0$. Given also that $f(d_0)=0$, we have $f(2d_0)=f(3d_0)=-d_0$, contradicting injectivity. 
    	\end{itemize}
    	\end{proof}
        To see why this claim uniquely determines $f$, we construct each $f(1), f(2), \cdots$ as follows: 
        for each $n$, assume that $f(d)=-d$ for all $d=1, 2, \cdots, n - 1$, 
        the injectivity claim implies that the only number smaller than $-(n-1)$ and is a divior of $n$ is $-n$. 
    	
    	\item [N5.]
    	The sequence of triples of positive integers $(a_0,b_0,c_0),(a_1,b_1,c_1),\ldots,(a_{2025},b_{2025},c_{2025})$ satisfies
    	$$(a_{i+1},b_{i+1},c_{i+1})=\big(a_{i}+ \gcd\left(b_{i},c_{i}\right),b_{i}+ \gcd\left(c_{i},a_{i}\right),c_{i}+ \gcd\left(a_{i},b_{i}\right)\big)$$for $0 \leqslant i \leqslant 2024.$ Determine the smallest possible value of $a_{2025}.$
    	
    	\textbf{Answer.} $3\cdot 2^{675}$. 
    	
    	\textbf{Solution.}
    	\textit{Construction}: take $(a_0, b_0, c_0)=(3, 5, 6)$, then 
    	\[
    	(3, 5, 6) \to (4, 8, 7) \to (5, 9, 11) \to (6, 10, 12)
    	\]
    	i.e. doubling every 3 sequences. 
    	
    	\textit{Bound}: note the map $(a_n, b_n, c_n) \to (a_{n+1}, b_{n+1}, c_{n+1})$ is permutation equivariant and scale preserving 
    	(i.e. if $(ka_n, kb_n, kc_n)\to (ka_{n+1}, kb_{n+1}, kc_{n+1})$. 
    	In modulo 2, we have the following iterations: 
    	\[
    	(0, 1, 1) \to (1, 0, 0) \to (1, 1, 1) \to (0, 0, 0)
    	\]
    	Each triple with $(a, b, c)$ not all even falls into one of the first three (up to permutations), i.e. $g(a, b, c) :=\min \{\nu_2(a), \nu_2(b), \nu_2(c)\}$ (the highest power of 2 dividing all entries) increases in at most 3 iterations. 
    	Thus for $i=0, 1, 2$ we have $a_{3k+(i+1)}-a_{3k+i}\ge 2^k$ as $g(a_{3k+i}, b_{3k+i}, c_{3k+i})\ge k$. 
    	Thus we have 
    	\[
    	a_{2025}-a_0\ge \sum_{i=0}^{674} (2^i+2^i+2^i)
    	=3(2^{675}-1)
    	\]
    	and given also $2^{675}\mid a_{2025}$ we have $a_{2025}\ge 2^{675}$. 
    	
    	
    	\item [N6.]
    	For positive integers $n$ and $k$, let $f_{k}(n)$ denote the remainder when $n$ is divided by $2^{k}-1.$ A positive integer $n$ is grouseif
    	\[ f_{1}(n)\le f_{2}(n)\le f_{3}(n)\le f_{4}(n)\le\dots \]
    	
    	\begin{enumerate}
    		\item [(a)]  Prove that there are infinitely many grouse numbers.
    		
    		\item [(b)] Prove that there exists a positive integer $M$ such that there are at most $\frac{M}{2025^{2025}}$ grouse numbers less than $M$.
    	\end{enumerate}
    	
    	\textbf{Solution.} 
    	(a) 
    	For each $k\ge 2$, consider the following number obtained by the following protocol: 
    	\[
    	s_0 = 2^{k}-1
    	\qquad 
    	s_j = \lceil\frac{s_{j-1}}{2^{k - j} - 1} \rceil \cdot (2^{k - j} - 1), 
    	\forall j = 1, 2, \cdots, k - 1
    	\]
    	We claim that $s_{k-1}$ is grouse, 
    	and that $2^k - 1\le s_{k-1} < 2^{k+1}-1$. 
    	Indeed, we have $0\le s_j - s_{j-1}\le 2^{k - j} - 2$. 
    	This means that $0\le s_{k - 1} - s_j = \sum_{i=1}^{k - j - 1} (2^{i} - 2) = 2^{k - j} - 2(k - j) < 2^{k - j} - 1$. 
    	Given also that $s_j$ is divisible by $2^{k - j} - 1$, 
    	we have $f_{k - j}(s_{k-1}) = s_{k - 1} - s_j$. 
    	Consequently, given that $s_0\le s_1\le \cdots \le s_{k - 2}\le s_{k-1}$, 
    	we have $f_1(s_{k-1})\le f_2(s_{k-1})\le \cdots \le f_{k-1}(s_{k-1})\le f_k(s_{k-1})$. 
    	Finally, given that $s_{k-1}\le s_0 + 2^k - 2k \le 2^{k+1} - 2k < 2^{k+1}-1$, 
    	we have $f_m(s_{k-1}) = s_{k-1}$ for $m\ge k + 1$. 
    	Thus this $s_{k-1}$ is grouse and there is at least one grouse number between the range $2^k - 1\le n < 2^{k+1}-1$ for all $k\ge 2$. 
    	
    	 (b) We show that if $n\ge 2^k - 1$ is grouse, then $n\ge (2^{m_1} - 1)(2^{k + 1 - m_1}-1)$ for all integers $m_1=1, 2, \cdots, k$. 
    	 A quick observation is the following. 
    	 
    	 \textit{If a number $s$ is grouse and $s\ge c\cdot (2^m-1)$ then $s\ge t$ with $t=(2^{m-1}-1)\cdot \lceil c\frac{2^m-1}{2^{m-1}-1}\rceil$}. 
    	 Indeed, if $s = c\cdot (2^m-1) + u$ where $u\ge 0$ then $f_m(s)\le u$ but if $s < t$, 
    	 then $f_{m-1}(c\cdot (2^m-1)) > 0$ and 
    	 then $f_{m-1}(s) = f_{m-1}(c\cdot (2^m-1)) + u > u$ so this $s$ is not grouse. 
    	 
    	 Plugging this back, we may do induction with $m_1$ (starting from $m_1=1$) and the conclusion follows from that 
    	 $\lceil (2^{m_1}-1)\frac{2^{k+1-m_1}-1}{2^{k-m_1}-1}\rceil\ge 1 + 2(2^{m_1}-1)=2^{m_1+1}-1$. 
    	 
    	 To finish, this means for any $m_1$ there are at most $2^{m_1} + 2^{k+1-m_1}$ grouse numbers in the range $\{2^k-1, \cdots, 2^{k+1}-2\}$. 
    	 Choosing $m_1=\lceil \frac{k+1}{2}\rceil$, we deduce that there is a constant $c > 0$ such that for any $M > 0$ there are at most $c\sqrt{M}$ grouse numbers that is at most $M$. 
    	 This solves the problem. 
    	
    	\item [N7.]
    	Let $\mathbb{N}$ denote the set of positive integers. A function $f\colon\mathbb{N}\to\mathbb{N}$ is said to be bonza if
    	\[
    	f(a)~~\text{divides}~~b^a-f(b)^{f(a)}
    	\]for all positive integers $a$ and $b$.
    	
    	Determine the smallest real constant $c$ such that $f(n)\leqslant cn$ for all bonza functions $f$ and all positive integers $n$.
    	
    	\textbf{Answer.} $c = 4$. 
    	
    	\textbf{Solution.}
    	\textit{Bound.} Set $a = b$, then $f(a)\mid a^a$. Suppose for some prime $p$ we have $p\mid f(a)$ (which then means $p\mid a$). 
    	Then take any $b$ divisible by $p$, $p\mid f(a)\mid b^a - f(b)^{f(a)}$ so $p\mid f(b)$. 
    	This means given a prime $p$, if $p\mid f(a)$ for some $a$ then $p\mid f(a)$ for all $a$ with $p\mid a$. 
    	Let $S$ be the set of primes dividing at least one value $f(a)$. 
    	Crucially, we note that for $p\in S$, $p\mid f(p)\mid p^p$ for any $p\in S$ so $f(p)$ is a (positive) power of $p$. 
    	Note that $p^k\equiv 1 \pmod{p}$, so $b^{p^k}\equiv b\pmod{p}$ by Fermat's little theorem, giving the following: 
    	\[
    	p\mid f(p) \mid b^p - f(b)^{f(p)}\equiv b - f(b)\pmod{p}
    	\]
    	for any $b$, so $p\mid b - f(b)$. 
    	
    	Assume $S$ is infinite. 
    	Then for each $b$, and consider $p\in S$ that's greater than both $b$ and $f(b)$. 
    	The fact $p\mid b - f(b)$ forces $f(b)=b$, hence $f$ is an identity function and so $c = 1$ in this case. 
    	
    	In another case, let $S$ be finite. Note that $S=\emptyset$ implies $c = 1$ too since $f\equiv 1$, 
    	so assume $S$ is nonempty. Let $p_0$ be the largest element of $S$. 
    	Then for all primes $q > p_0$ we have $f(q)=1$ so $p_0\mid q - 1$ for all such primes $q$. 
    	However, by Dirichlet's theorem, for any $r$ with $\gcd(r, p_0)=1$ there are infinitely many primes $q$ such that $p_0\mid q - r$, 
    	so $\gcd(r, p_0)=1\to p_0 \mid r - 1$. This forces $p_0 = 2$. 
    	In other words, $f(a)$ is a power of two for all $a$, with $f(a)=1$ for all odd numbers $a$. 
    	Now set $b = 3$ and let $a = c\cdot 2^k$ be even, with $c = 1$. Then 
    	\[
    	f(a) \mid 3^a - 1 = (3^c - 1)(3^c+1)\cdots (3^{2^{k-1}c}+1)
    	\]
    	and note that $3^c-1, 3^{2c}+1, \cdots, 3^{2^{k-1}c}+1$ are all divisible by 2 but not by 4, and $3^c+1$ divisible by 4 but not by 8. 
    	Thus $\nu_2(3^a-1)=k+2$, i.e. $f(a)\le 2^{k+2}\le 4a$. 
    	Given also $f(a)=1$ for all odd $a$ we have $c\le 4$. 
    	
    	\textit{Construction.} The example above motivates the following construction: 
    	\[
    	f(a)=
    	\begin{cases}
    		1 & a\text{ odd}\\
    		4 & a=4k+2\\
    		16 & 4\mid a\\
    	\end{cases}
    	\]
    	Here $f(4)=16$ suggests $c = 4$ cannot be improved. To verify that this function works, we need to show that 
    	\[
    	4\mid b^a-f(b)^{4}, a = 4k+2
    	\qquad 
    	16\mid b^a - f(b)^{16}, a = 4k
    	\
    	\]
    	In the first case we automatically have $4\mid b^a, f(b)^4$ when $a, b$ both even, 
    	and when $b$ is odd we have $b^a\equiv 1\pmod{4}$ (since $b^a$ is a square). 
    	In the second case we have $16\mid b^a, f(b)^4$ when $b$ even and $a\ge 4$, 
    	and when $b$ is odd we note that $b^4-1=(b-1)(b+1)(b^2+1)$ is divisible by 16. 
    	
    	\item [N8.]
    	Prove that there are finitely many prime numbers $p$ such that
    	
    	\begin{itemize}
    		\item $p-8$ is a perfect square, and
    		
    		\item for all odd positive integers $k<\sqrt{p}$, the number $p-k^2$ has at most two distinct prime factors.
    	\end{itemize}
    	
    	\textbf{Solution.}
    	We consider the following claim. 
    	
    	Claim: all such $p$'s must be in the form of $4\cdot 2^k\cdot (2^k\pm 1) + 9$ where the last factor $2^k\pm 1$ is odd prime. 
    	
    	Proof: note that $p$ is odd and $p - 9 = a^2 - 1$ for some $a$ (hence even), 
    	so can have at most one odd factor. 
    	Given $\gcd(a-1, a+1)=2$, one of them is power of 2 and the other one is $2\cdot$ odd prime power. 
    	But the only two consecutive perfect powers are 8 and 9, 
    	hence this `prime power' must be itself a prime. 
    	
    	Now we may split into the following cases. 
    	
    	Case $2^k + 1$: for this to be prime, $k$ must be power of 2, i.e. $k = 2^m$. 
    	Then $p = 4\cdot 2^k (2^k + 1) + 9 = (2^{k + 1} + 1)^2 + 8$. 
    	When considering modulo 17, 
    	for $k = 2^m$ with $m\ge 3$ we have 
    	\[
    	(2^{2^m+1} + 1)^2 + 8\equiv (2 + 1)^2+8\equiv 0\pmod{17}
    	\]
    	as $2^8\equiv 1\pmod{17}$, thus this $p$ is not prime for all such $p$'s. 
    	
    	Case $2^k - 1$: for this to be prime, $k$ must be itself a prime (in particular, $k$ is odd). 
    	We have a similar expansion $p = (2^{k+1}-1)^2+8$. 
    	
    	Claim: when $k\ge 7$, there exists a prime number $q\le 2p^{1/4}$ such that $p$ is a quadratic residue modulo $p$. 
    	
    	Proof: we will find those $q$ where $q\mid 2^{k+1}-1$, and 2 (and therefore 8) is a quadratic residue. 
    	When $k + 1$ is not a power of 2, let $m > 1$ be any odd divisor of $k + 1$, then $2^m-1 \mid 2^{k+1}-1$. 
    	Take any prime $q\mid 2^m - 1$, since $m$ is odd, we have order modulo $\text{ord}_q(2)$ odd as $\text{ord}_q(2)\mid m$. 
    	But $\phi(q)=q-1$ is even, so 2 is indeed a quadratic residue. 
    	
    	When $k+1$ is a power of 2, 
    	we have the expansion $2^{k+1}-1 = (2^{\frac{k+1}{2}} - 1)(2^{\frac{k+1}{2}} + 1)$ and $2^{\frac{k+1}{2}} + 1$. 
    	We now use the identity that a) 2 is a quadratic residue for prime $q$ if and only if $q\equiv 1, 7\pmod{8}$, 
    	and b) all prime divisors of $x^{2^m}+1$ are of the form $\ell \cdot 2^{m+1} + 1$, 
    	as $\text{ord}_q (x)=2^{m+1}$. 
    	Thus when $k+1\ge 8$ (i.e. $m\ge 2$), 2 is a quadratic residue of any such prime $q$. 
    	
    	Taken together, all such $q$'s are at most $2^{\frac{k+1}{2}} + 1\le 2(2^{\frac{k+1}{2}} - 1)\le 2\sqrt{2^{k + 1} - 1}\le 2p^{1/4}$, 
    	proving our claim. 
    	
    	
    \end{enumerate}
\end{document}